\documentclass[11pt]{article}
\usepackage[margin=1in]{geometry}
\usepackage[T1]{fontenc}
\usepackage[utf8]{inputenc}
\usepackage{lmodern}
\usepackage{microtype}
\usepackage{booktabs}
\usepackage{array}
\usepackage{enumitem}
\usepackage{amsmath}
\usepackage{hyperref}
\usepackage{xcolor}

\hypersetup{
  colorlinks=true,
  linkcolor=blue!60!black,
  urlcolor=blue!70!black,
  citecolor=blue!60!black
}

\setlist[itemize]{leftmargin=1.2em}
\setlist[enumerate]{leftmargin=1.4em}

\title{Group 13 Report:\\Turn-Level User Intent Annotation (MultiWOZ 2.2)}
\author{
  Mayowa Adesanya (\texttt{adesanym}, 400233110) \and
  Nicholas Zajkeskovic (\texttt{zajkeskn}, 400402381) \and
  Divij Dhiraaj (\texttt{dhiraajd}, 400360225)
}
\date{\today}

\begin{document}
\maketitle

\section{Group Number and Members}
\textbf{Group Number:} 13
\begin{itemize}
  \item Mayowa Adesanya (\texttt{adesanym}, 400233110)
  \item Nicholas Zajkeskovic (\texttt{zajkeskn}, 400402381)
  \item Divij Dhiraaj (\texttt{dhiraajd}, 400360225)
\end{itemize}

\section{Annotation Instructions}
\subsection{Task Definition}
For each instance, annotators read:
\begin{enumerate}
  \item \texttt{system\_context} (previous system turn)
  \item \texttt{user\_utterance} (current user turn)
\end{enumerate}
Annotators assign exactly one dominant intent label for the current user turn.

\subsection{Label Set}
\begin{enumerate}
  \item \texttt{REQUEST}: ask for information/options or request an action.
  \item \texttt{INFORM\_CONSTRAINT}: provide preferences/details (price, area, time, people, etc.).
  \item \texttt{CONFIRM\_ACCEPT}: explicit acceptance/confirmation.
  \item \texttt{CORRECT\_CLARIFY}: correction or clarification after mismatch.
  \item \texttt{SOCIAL}: greeting/thanks/closing with no task-directed intent.
\end{enumerate}

\subsection{Operational Procedure}
Each annotator used one workbook:
\begin{itemize}
  \item \texttt{annotation/excel/adesanym\_annotation.xlsx}
  \item \texttt{annotation/excel/zajkeskn\_annotation.xlsx}
  \item \texttt{annotation/excel/dhiraajd\_annotation.xlsx}
\end{itemize}

Each workbook contains:
\begin{itemize}
  \item \texttt{initial} sheet (primary labels)
  \item \texttt{reannotation} sheet (overlap labels)
\end{itemize}

Per-row process:
\begin{enumerate}
  \item Read context and utterance.
  \item Apply decision order.
  \item Choose one label from dropdown.
  \item Add concise note only when ambiguous.
  \item Do not edit metadata columns.
\end{enumerate}

\subsection{Decision Order and Tie-Break Rules}
\textbf{Decision order:}
\begin{enumerate}
  \item \texttt{CORRECT\_CLARIFY}
  \item \texttt{CONFIRM\_ACCEPT}
  \item \texttt{REQUEST}
  \item \texttt{INFORM\_CONSTRAINT}
  \item \texttt{SOCIAL}
\end{enumerate}

\textbf{Tie-break policy (mixed intents):}
\begin{enumerate}
  \item Correction beats all other intents.
  \item Explicit acceptance beats request/inform.
  \item Request beats inform when asking for next system action.
  \item Inform beats social when task content exists.
  \item Social only when no task-directed intent remains.
\end{enumerate}

\section{Dataset Description and Sensitive-Content Disclaimer}
\subsection{Source and Provenance}
The dataset was built from MultiWOZ 2.2 dialogues.
\begin{itemize}
  \item Dataset GitHub: \url{https://github.com/budzianowski/multiwoz}
  \item Raw hosted files used in downloader:
  \url{https://raw.githubusercontent.com/budzianowski/multiwoz/master/data/MultiWOZ_2.2}
\end{itemize}

\subsection{Annotation Unit}
Each instance is one pair:
\begin{itemize}
  \item previous system utterance (\texttt{system\_context})
  \item current user utterance (\texttt{user\_utterance})
\end{itemize}

\subsection{Annotated Data Volume}
\textbf{Annotated dataset submitted:} \texttt{data/processed/annotations\_labeled\_long.csv}
\par\medskip

\begin{center}
\begin{tabular}{lrr}
\toprule
\textbf{Component} & \textbf{Count} & \textbf{Notes} \\
\midrule
Initial (unique) annotations & 1050 & 350 per annotator \\
Reannotation overlap entries & 135 & 45 per annotator \\
Total label entries & 1185 & initial + overlap \\
Blank labels after completion & 0 & quality check passed \\
\bottomrule
\end{tabular}
\end{center}

\subsection{Sensitive-Content Disclaimer}
Although primarily benign, the dialogue data may contain rude or impolite phrasing, abrupt tone shifts, and occasional potentially offensive language. The corpus is used strictly for academic annotation and analysis.

\section{Interesting Data Points, Estimates, and Agreement}
\subsection{Required Annotation Estimates}
\begin{itemize}
  \item Total unique instances annotated: 1050
  \item Total labeling operations including overlap: 1185
\end{itemize}

Using the agreed 1-hour baseline per annotator for the initial pass:
\[
\text{Average time per instance}=\frac{60 \times 60}{350}=10.29\text{ seconds}
\]

\subsection{Agreement Metric Selection}
Because overlap was distributed among three annotators (rather than a single fixed annotator pair), we used \textbf{Krippendorff's Alpha (nominal)} on overlap instances.

\begin{center}
\begin{tabular}{lr}
\toprule
\textbf{Metric} & \textbf{Value} \\
\midrule
Overlap instances & 135 \\
Percent agreement & 0.7037 \\
Krippendorff's Alpha (nominal) & 0.5997 \\
\bottomrule
\end{tabular}
\end{center}

Interpretation: reliability is \textbf{moderate}. The labels are usable for analysis, but disagreement remains meaningful in boundary cases.

\subsection{Three Interesting Data Points}
\subsubsection*{Example 1: Hostility + task redirection}
\textbf{Instance:} \texttt{mw22\_PMUL4392.json\_2}\\
\textbf{User turn:} ``uhm, no. that's your job, weirdo. also, i need a train.''\\
\textbf{Labels:} Original \texttt{REQUEST}, Reannotation \texttt{CORRECT\_CLARIFY}.\\
\textbf{Why interesting:} combines rejection/correction with a new request, creating dominant-intent ambiguity.

\subsubsection*{Example 2: ``instead'' ambiguity}
\textbf{Instance:} \texttt{mw22\_MUL1543.json\_6}\\
\textbf{User turn:} ``Can you help me find a train going to Stevenage and leaving Thursday instead?''\\
\textbf{Labels:} Original \texttt{REQUEST}, Reannotation \texttt{CORRECT\_CLARIFY}.\\
\textbf{Why interesting:} ``instead'' can indicate correction while turn form remains request-like.

\subsubsection*{Example 3: Social preface + new intent}
\textbf{Instance:} \texttt{mw22\_PMUL0719.json\_12}\\
\textbf{User turn:} ``Thank you. I also want a train for friday, from broxbourne to cambridge.''\\
\textbf{Labels:} Original \texttt{INFORM\_CONSTRAINT}, Reannotation \texttt{REQUEST}.\\
\textbf{Why interesting:} social phrase and actionable intent co-occur in one turn.

\subsection{Disagreement Patterns}
\textbf{Most frequent disagreement transitions:}
\begin{itemize}
  \item \texttt{REQUEST} $\rightarrow$ \texttt{CONFIRM\_ACCEPT} (8)
  \item \texttt{INFORM\_CONSTRAINT} $\rightarrow$ \texttt{REQUEST} (6)
  \item \texttt{REQUEST} $\rightarrow$ \texttt{CORRECT\_CLARIFY} (4)
\end{itemize}

\textbf{Per-original-label agreement:}
\begin{center}
\begin{tabular}{lr}
\toprule
\textbf{Original Label} & \textbf{Agreement Rate} \\
\midrule
\texttt{REQUEST} & 72.7\% \\
\texttt{INFORM\_CONSTRAINT} & 69.0\% \\
\texttt{CONFIRM\_ACCEPT} & 52.6\% \\
\texttt{CORRECT\_CLARIFY} & 25.0\% \\
\texttt{SOCIAL} & 95.8\% \\
\bottomrule
\end{tabular}
\end{center}

\section{Reflection Questions}
\subsection*{(a) What did you learn about the task from doing the annotation?}
We learned that intent annotation depends strongly on conversational context, not only keywords. Short forms such as ``yes'' and ``thanks'' require prior-turn grounding. Deterministic tie-break rules reduced inconsistency in mixed-intent turns.

\subsection*{(b) What challenges do you expect models to face when learning from your data?}
Key challenges are mixed-intent turns, informal language variation, context dependence for short responses, and boundary overlap between \texttt{REQUEST}, \texttt{INFORM\_CONSTRAINT}, and \texttt{CORRECT\_CLARIFY}. The relatively small count of \texttt{CORRECT\_CLARIFY} also increases learning difficulty.

\subsection*{(c) What surprising things did you observe in the data?}
We observed high consistency on clearly social turns, lower consistency on correction-heavy turns, and frequent pragmatic shifts where a single utterance both responds to prior context and introduces a new goal.

\subsection*{(d) Which features do you expect to be useful?}
Useful features likely include lexical action cues (book/find/need/instead/no), confirmation and rejection markers, explicit slot constraints (time/day/location/people), prior-turn context features, and dialogue-level intent transition patterns.

\subsection*{(e) Describe the mental model you used for annotation.}
We used a decision ladder:
\begin{enumerate}
  \item detect correction first,
  \item then explicit acceptance,
  \item then actionable request,
  \item then constraint-provision,
  \item else social.
\end{enumerate}
Ambiguous cases were resolved by tie-break order and documented in notes.

\subsection*{(f) Was anything unclear about the instructions, and how would you improve them?}
Initial high-level instructions were usable but not fully self-contained for edge cases. The most helpful refinements were explicit decision ordering, explicit tie-break rules, compact examples for mixed intents, and explicit limits on editable spreadsheet fields.

\section{Conclusion}
The group completed annotation and overlap labeling, then computed agreement on the overlap set. The resulting Krippendorff's Alpha of 0.5997 indicates moderate reliability: strong for some classes (especially social turns) and weaker for correction/request boundaries. The dataset is suitable for baseline model training and error analysis, with clear guidance for future annotation guideline refinement.

\end{document}
